\documentclass[10pt,a4paper]{amsart}
\usepackage[margin=19mm]{geometry}
\usepackage{amsmath,amssymb,mathtools}
\usepackage[scr=rsfs,cal=euler]{mathalfa}
\usepackage{xfrac}
\usepackage[unicode,hidelinks,bookmarks=false]{hyperref}
\newcommand{\ie}{i.e.\ }
\newcommand{\cf}{cf.\ }
\newcommand{\resp}{respectively\ }
\newcommand{\Subsection}{\subsection}
\providecommand{\nobreakdash}{\nobreak\hbox{-}}
\setlength{\emergencystretch}{2em}
\multlinegap0pt
\usepackage{amssymb}
\usepackage{algorithm}\usepackage{algpseudocode}
\usepackage{bm}
\newtheorem{thm}{Theorem}
\newcommand{\CC}{\mathbf{C}} % complex numbers
\def\GL{\text{GL}}
\newcommand{\KK}{K}
\newcommand{\OK}{{{\mathcal O}_K}}
\providecommand{\binom}[2]{{#1\choose#2}}
\newcommand{\bth}{{{\theta_*}}} %or $\Theta$,
\newcommand{\calb}{\mathcal{B}}
\newcommand{\cals}{\mathcal{S}}
\renewcommand{\geq}{\geqslant}
\renewcommand{\leq}{\leqslant}
\newcommand{\set}[1]{\{#1\}}
\title{When are Permutation Invariants Cohen-Macaulay?}
\author{H.E.A. Campbell, David L. Wehlau}
\date{Source: arXiv:2308.09056v2 (CC BY 4.0)}
\begin{document}
\maketitle

\noindent\emph{Source and licence.} When are Permutation Invariants Cohen-Macaulay?. Version arXiv:2308.09056v2 (CC BY 4.0), distributed by arXiv under the Creative Commons Attribution 4.0 International licence, http://creativecommons.org/licenses/by/4.0/, as recorded in the arXiv licence metadata for this exact version and shown on the abstract page https://arxiv.org/abs/2308.09056v2. Full licence notice: \texttt{licenses/arxiv-2308.09056-CC-BY-4.0.txt}.\par\smallskip

\noindent\textit{Editorial scope.} Counted unit: the article's thm thm (source0.tex lines 1104--1106). The article's complete original proof of it is delivered with it (source0.tex lines 1108--1157, one \texttt{proof} environment of 50 lines). No mathematics is written, repaired or re-derived: the statement and the proof are the article's own lines, in the article's own notation, and no retained line is substituted or re-flowed.  No further source-local context is needed: every cross-reference of every retained line resolves inside the two delivered spans.  Nothing was omitted from any retained span, so the package carries no omission record.  No retained line contains a citation, so the wrapper carries no bibliography and no bibliography entry.  No retained line contains a figure environment, an image-inclusion command or drawing code, so no image is delivered and the package declares no asset dependency.  Editorial formatting only, in addition: an AMS \texttt{amsart} preamble that restates the article's own declarations and macros used by the retained lines, namely the environments \texttt{thm} on separate counters, together with the standard packages those lines need.  The retained environments print in this document as Theorem 1 in order of appearance, whereas the article prints the same items with the numbering of its own full document.  The complete native source of the article as distributed in the arXiv e-print for this exact version is bundled unchanged as \texttt{sources/145774/source0.tex}.
\begin{thm}
$\lambda_\calb^{e(G,\calb)}$ divides $\det M(\bth)$ in $\KK[V]$.  
\end{thm}

 \begin{proof}
 Note that $\Sigma$ acts on the set of rows of $M(\theta_*)$ by permutations.
 Let $W$ denote the $\CC$-vector space spanned by the rows of $M(\theta_*)$.  
 This is a subspace of $\oplus_{i=1}^\ell \CC[V]_{a_i}$.  
 We examine $W$ as a permutation representation of the cyclic group $C_P$ of order $c$.
 Since $C_P$ is abelian there is a basis $\set{r_1,r_2,\dots,r_\ell}$ of $W$ consisting of eigenvectors for $g_P$,
 i.e., working over $\CC$, we may diagonalize the action of $g_P$  on $W$.   Then $g_P r_i = \mu_i(g_P) r_i$ for some character $\mu_i$ of $C_P$.  We fix a choice of primitive $c^{\rm{th}}$ root of unity $\zeta_c\in\CC$ by writing $\det(g_P)=\zeta_c^{-1}$. 
 Note that since $g_p \in \GL(n,\OK)$, we have that $\zeta_c \in \OK$.

 The change of basis matrix $U$ which diagonalizes the 
 action of $g_P$ on W is an
invertible $\ell \times \ell$ matrix with entries in $K$.
Writing $\widetilde M := U M(\bth)$, we see that 
$\det (\widetilde M) = \kappa \det(M(\bth))$ where 
$\kappa = \det(U) \in K$.

We extend the set $\{L_P\}$
to a basis $\{L_P,L_2,\dots,L_n\}$ of $V^*$ where $\{L_2,\dots,L_n\}$ is dual to a basis for the reflecting hyperplane $P$.
Then $g_P \cdot L_P = \zeta_c L_P$ and $g_P \cdot L_i = L_i$
for $2 \leq i \leq n$.  
In this basis, every monomial is an eigenvector for $C_P$
with eigenvalue given by $\zeta_c^k$ with 
$0 \leq k \leq c-1$ where the degree of the monomial in $L_P$
is congruent to $k$ mod $c$.



 Now the cycle structure of $g_P$ acting on $\cals$ determines the permutation action of $C_P$ on $W$.
 A cycle of length $a$ corresponds to a transitive $a$ dimensional permutation $C_P$-subrepresentation of $W$.
 
  Diagonalizing the action of $g_P$ on this $a$ dimensional subrepresentation reveals eigenvectors for $g_P$
 with eigenvalues: $\zeta_{c}^{c/a}, \zeta_{c}^{2c/a},\dots,\zeta_{c}^{ac/a}=1$.
Thus the eigenvalue associated to $r_i$ is
$\mu_i(g_P)=\zeta_c^{cj/a}$ for some $j$ with
$0 \leq j < a$.  
In particular every monomial (when written in the basis $L_P, L_2,\dots, L_n$) appearing in 
each of the $\ell$ entries in $r_i$ must be divisible by $L_P^{c j/a}$.

Therefore a transitive permutation subrepresentation of $W$ of dimension $a$ 
contributes the factor 
$L_P^{\sum_{i=0}^{a-1} c i/a} = L_P^{c\binom{a}{2}/a}=L_P^{c(a-1)/2}$ to $\det \widetilde{M}$. 
 Taking all the cycles into account we see that
 $L_P^{(c/2)\sum_{i \geq 2} b_i(i-1)}$ divides 
 $\det \widetilde{M}$.
 Since for any two hyperplanes $P_1$ and $P_2$, the linear forms
 $L_{P_1}$ and $L_{P_2}$ are co-prime,  
  $\lambda_\calb^{e(G,\calb)}$ divides $\det \widetilde{M}$ in $\KK[V]$ for all orbits $\calb$.
  Therefore $\lambda_\calb^{e(G,\calb)}$ divides 
  $\det M(\bth)$ in $\KK[V]$ for all orbits $\calb$.
\end{proof}

\end{document}
